% Generated draft: requires mathematical and source review.
\documentclass[11pt]{article}
\usepackage[margin=25mm]{geometry}
\usepackage{fontspec}
\setmainfont{Latin Modern Roman}
\usepackage{amsmath,amssymb,mathtools,unicode-math}
\setmathfont{Latin Modern Math}
\usepackage{enumitem,needspace,xurl,hyperref}
\setlength{\parindent}{0pt}
\setlength{\parskip}{5pt}
\setlength{\emergencystretch}{3em}
\newcommand{\N}{\mathbb{N}}
\newcommand{\Z}{\mathbb{Z}}
\newcommand{\Q}{\mathbb{Q}}
\newcommand{\R}{\mathbb{R}}
\newcommand{\C}{\mathbb{C}}
\newcommand{\F}{\mathbb{F}}
\newcommand{\A}{\mathbb{A}}
\newcommand{\PP}{\mathbb{P}}
\newcommand{\E}{\mathbb{E}}
\newcommand{\eps}{\varepsilon}
\newcommand{\dd}{\,\mathrm{d}}
\begin{document}
\section*{UnsolvedMath ID: 11000299}
\textbf{Source problem code:} catalogue entry
\subsection{Definitions and mathematical statement}
Abelianization \(F_n\to\Z^n\) induces
\[
 \operatorname{Out}(F_n)\longrightarrow\operatorname{GL}(n,\Z).
\]
Define \(IA_{\mathrm{src}}(n)\) to be its kernel, exactly as in the source. The principal question is whether \(IA_{\mathrm{src}}(n)\) has a finite presentation for every \(n>3\), and more broadly to establish its finiteness properties. A finite presentation means finitely many ordinary generators and defining relations. More generally, type \(F_k\) means that a \(K(IA_{\mathrm{src}}(n),1)\) has finite \(k\)-skeleton; \(F_1\) is finite generation and \(F_2\) finite presentability.

The notation must not be confused with the often-denoted \(IA_n\) kernel in \(\operatorname{Aut}(F_n)\). For \(n\geq2\), the latter fits into
\[
 1\longrightarrow\operatorname{Inn}(F_n)\cong F_n
 \longrightarrow\ker(\operatorname{Aut}(F_n)\to\operatorname{GL}(n,\Z))
 \longrightarrow IA_{\mathrm{src}}(n)\longrightarrow1.
\]
The source compares its outer kernel to a surface Torelli group but supplies no further finiteness theorem.

The research summary asserts finite generation and finite presentation for \(n\geq3\), and leaves higher \(F_k\) open, naming several authors without exact papers. These claims do not identify which kernel or presentation notion is used. In particular, finite invariant/L-presentations or finite relation orbits are not the same as ordinary finite presentations. The summary is therefore retained as an unverified, potentially conflated progress report rather than used to resolve the printed \(n>3\) question.
\subsection{Short English statement}
Determine finiteness properties, especially ordinary finite presentability for \(n>3\), of the outer-automorphism Torelli kernel, keeping it distinct from the automorphism-group kernel.
\subsection{Sources}
\begin{itemize}
\item \emph{Problems on Mapping Class Groups and Related Topics}, edited by Farb, Question 6.3, supplied PDF page 338. \url{https://www.math.uchicago.edu/~farb/papers/mcgbook.pdf}. Recorded access 2026-07-29; extraction pdf-font-structure.
\item The supplied summary names Magnus, Cohen--Pakianathan and Andreadakis for generation-related work, and Day--Putman and Wang for a claimed finite-presentation result, without complete citations or exact theorem/presentation notions.
\item Local UnsolvedMath registry, \texttt{data/unsolvedmath/problems.json}, permanent ID \(11000299\), exact source code \texttt{catalogue entry}; registry OPEN-TRIAGE, with a PARTIAL-PROGRESS summary asserting \(F_2\) known.
\end{itemize}
\end{document}
